\section{Rencana Kegiatan}
Rencana kegitana yang akan saya lakukan adalah sebagia berikut:
\begin{itemize}
    \item \textbf {Studi literatur} \\ Peneliti akan melakukan kajian terhadap berbagai literatur dan penelitian terdahulu yang relevan dengan prediksi curah hujan,  \textit{spatial weight matrix} serta penggunaan algoritma STARMA. Studi ini bertujuan untuk memahami landasan teori yang telah dikembangkan serta mengevaluasi kelebihan dan kekurangan pendekatan yang digunakan pada penelitian sebelumnya. Referensi utama akan mencakup jurnal internasional dan publikasi ilmiah
    \item \textbf {Pengumpulan Data} \\ Data yang digunakan adalah data curah hujan bulanan bersifat spasial dan temporal yang diperoleh dari sumber data sekunder seperti NASA POWER. Data tersebut akan mencakup beberapa titik lokasi (per Kecamatan di empat bagian Surabaya) dalam kurun waktu 10 tahun terakhir, sehingga memungkinkan dilakukannya analisis spatio-temporal secara menyeluruh. 
    \item \textbf{Perancangan Modeling} \\ Pada tahapan ini, peneliti akan melakukan eksplorasi dan analisis data, termasuk visualisasi data spasio-temporal, pengujian kestasioneran, serta perancangan matriks bobot spasial \textit{(spatial weight matrix)}. Selanjutnya, dilakukan desain model prediktif berbasis STARMA yang menggabungkan komponen spasial dan temporal. 
    \item \textbf{Implementasi Modeling} \\ Implementasi model dilakukan menggunakan perangkat lunak seperti python dan R dengan pustaka pendukung untuk analisis spasio-temporal. Model STARMA yang telah dirancang akan diterapkan pada data curah hujan dan dilakukan pelatihan serta validasi model untuk mendapatkan parameter optimal. 
    \item \textbf{Analisa Hasil Implementasi} \\ Setelah model diterapkan, hasil prediksi curah hujan akan dianalisis berdasarkan tingkat akurasi, kesesuaian dengan pola historis serta kemampuan model dalam menangkap dinamika spasio-temporal.
    \item \textbf{Penulisan Laporan } \\ Seluruh proses penelitian mulai dari perancangan, implementasi hingga analisis hasil akan didokumentasikan secara sistematis dalam bentuk laporan akhir. Laporan akhir ini mencakup latar belakang, rumusan masalah, metodologi, hasil, pembahasan dan kesimpulan.
    \
Analisis dan perancangan sistem 
\end{itemize}